
\documentclass[12pt]{article}
\usepackage[UTF8]{ctex}
\usepackage[a4paper,margin=2.5cm]{geometry}
\usepackage{enumitem}
\pagestyle{plain}
\begin{document}
\begin{center}{\Large\bfseries 全国大学生数学建模竞赛\\[2mm] AI 工具使用详情说明}\end{center}
\section*{一、所用 AI 工具名称与版本}
DeepSeek 大模型驱动的数学建模智能体（DeepSeek Harness / CUMCM 专项预设 mathmodel-v7，模型：qwen3.8-flash）。
竞赛期间使用 AI 工具辅助，参赛队对作品的原创性、真实性、准确性负全部责任（按 2026 年参赛规则第 6 条与《人工智能工具使用规定》执行）。

\section*{二、具体使用目的和环节}
\begin{itemize}[leftmargin=2em]
\item 赛题文本与附件数据结构化解析（题面 PDF/Excel 读取登记，全部结果与官方原文逐字核对）；
\item 建模路线检索与竞争式筛选（多智能体独立路线生成、公平预算侦察、对抗评审），关键取舍均记录于决策日志并人工审查；
\item 代码辅助编写与调试（CP-SAT/校验器实现；所有数值结论均由本地真实运行产生，并经独立第二实现复核）；
\item 论文语言整理与 LaTeX 排版辅助（公式、图表引用由模型契约 FINAL\_MODEL\_SPEC 与验证计划逐项核对）。
\end{itemize}

\section*{三、主要提示方式与使用过程说明}
按阶段下达任务式提示（题面解析、建模设计、求解实现、验证、写作），每阶段产物由机器门禁（schema 校验、独立检查器、数值溯源）验证后方进入下一阶段；未向任何平台检索本届赛题题解或参赛讨论，检索仅限通用方法文献。

\section*{四、采纳、人工修改与核验情况}
\begin{itemize}[leftmargin=2em]
\item 关键决策示例（决策日志摘录，共 22 条）：
\item \textbf{[D-INTAKE-001]}（intake）：missing\_coverage(data) MEDIUM 人工确认放行→补 data/canonical\_plans.csv 后自动消除（现 0 冲突）
\item \textbf{[D-Q3-001]}（planned-analysis）：问题3 主口径：Q2 方案保持不动，新增 C 装备全视界自由落位（不限位置即题面 不限平移幅度）；新C装备参数沿用附件C模板（宽3/时长2/间隔8/次数12），另做口径灵敏度讨论
\item \textbf{[D-HORIZON-001]}（brainstorm-mathmodel）：时间视界裁定：T\_MAX=643Δt（slot 起点公式 t1+k*(g+d)，与题面 A001 例证一致）。DATASET\_SNAPSHOT.class\_profile.last\_slot\_end\_max=621 系公式漏 d 的陈旧字段，判定 stale：以附件1（primary\_official，哈希不变）运行时复算为准；脚本 code/intake\_dataset\_snapshot.py 已修正，coding 
\item \textbf{[D-BUDGET-001]}（brainstorm-mathmodel）：两级预算正式分离：①brainstorm 侦察预算 600s/候选（standard 档协议冻结件，仅用于路线间公平比较，其 incumbent/证书不作生产结果）；②coding\_visual 生产求解预算=长时后台运行（Q2/Q4 词典序分层+撤销阶梯计划 1-2 小时，允许分批续跑），最终数字只能来自生产运行并经独立 checker 复验
\item \textbf{[D-BUDGET-R51]}（brainstorm-mathmodel）：CH-REV-A 对 Q2-R51 预算越界申报的 Manager 裁定=有条件豁免：incumbent 求解 547.3s 在 600s 协议内；超出的 120.4s 为独立「界证书」进程（r=6 incumbent 未变、不替换解、仅收紧/开放界），属 bound 证据不属搜索预算。公平性成立（其余候选 incumbent 均 ≤600s 达成，超出部分未改善其 objective）。豁免登记于 IDEA\_DECISION；生产阶
\item \textbf{[D-R32-REVISE]}（brainstorm-mathmodel）：Q2-R32 verdict=revise 的处置：R32 不作 primary（被 R51 支配），revise 三项中「随机版补测」与「workers=1 复现」记为豁免（理由：其结论已被定位为负结果证据+热启动源，随机性不影响该定位）；「范式重定位」采纳——IDEA\_DECISION 中 R32 定位改写为『CP-SAT 主导混合范式，连续引导贡献量化为 297→175 的负结果证据』。R32 侦察结果保留为对比证据不作质量基准
\item \textbf{[D-PROD-REPRO]}（brainstorm-mathmodel）：评审遗留要求转生产验证：①Q2/Q4 生产 CP-SAT 运行必须做 workers=1 与多 worker 的 incumbent 一致性验证（复现性）；②撤销层 Σr≤5 判定在生产以 ≥1e4s 级预算或专用下界机制攻闭合，闭合前论文只写区间；③Q1 论文引用 R41 压力测试须措辞为『随机化召回诊断证据（n≈500，±4\%量级）』，R52 类内界称『错误探测界』。转 VALIDATION\_PLAN 条目。
\item \textbf{[D-Q3-WIN]}（brainstorm-mathmodel）：Q3 侦察关键结果：Q3-R11 集包装 CP-SAT 双模型 OPTIMAL incumbent=bound=139 + HiGHS LP 分数上界=139.0 + 候选放置 98×532 穷举完备（三路对拍 131254 干扰边 0 失配）⇒ 在 R5 主口径（Q2-R51 基座+模板参数+643 视界）下 N*=139 可证。R51 初等双计数界 332 合法但松 2.39×，仅作可手核对照。生产阶段将在新 Q2 生产解上重算（基
\item \textbf{[D-CLOSE-brainstorm]}（brainstorm-mathmodel）：阶段 brainstorm 收口（tx=SC-A71EA16E7E60）：outputs/gate/schema 全部通过
\item \textbf{[D-BS-CLOSE]}（brainstorm-mathmodel）：brainstorm 收口：winner Q1-R21/Q2-R51/Q3-R11/Q4-R11；等级 L2/L2/L3\_NEAR\_GLOBAL/L2。评审闭环完成：Q4-R11 锚定植入 Q2-R51 解（run\_anchor 四断言 ALL\_PASS，自跑8未支配植入6→生产热启动冻结收益判定）、Q4-R22 作为求解路线 rejected、Q2-R51 预算豁免 D-BUDGET-R51、R23 补 3-seed、R11 
\item \textbf{[D-CLOSE-analysis]}（2analysis-modeling）：阶段 analysis 收口（tx=SC-6797B89785F1）：outputs/gate/schema 全部通过
\item \textbf{[D-CLOSE-methodology_review]}（7methodology-review）：阶段 methodology\_review 收口（tx=SC-81920836DD07）：outputs/gate/schema 全部通过
\item \textbf{[D-CLOSE-methodology_review]}（7methodology-review）：阶段 methodology\_review 收口（tx=SC-494D6C15B5F2）：outputs/gate/schema 全部通过
\item \textbf{[D-CLOSE-validation_plan]}（8validation-plan）：阶段 validation\_plan 收口（tx=SC-12DEA893DA33）：outputs/gate/schema 全部通过
\item \textbf{[D-CLOSE-methodology_review]}（7methodology-review）：阶段 methodology\_review 收口（tx=SC-DF553EF5396C）：outputs/gate/schema 全部通过
\item \textbf{[D-CLOSE-validation_plan]}（8validation-plan）：阶段 validation\_plan 收口（tx=SC-154752A0AF92）：outputs/gate/schema 全部通过
\item \textbf{[D-CLOSE-validation_plan]}（8validation-plan）：阶段 validation\_plan 收口（tx=SC-3AFA1BF6CEDC）：outputs/gate/schema 全部通过
\item \textbf{[D-CLOSE-validation_plan]}（8validation-plan）：阶段 validation\_plan 收口（tx=SC-D85F20C330E3）：outputs/gate/schema 全部通过
\end{itemize}
\end{itemize}
\begin{itemize}[leftmargin=2em]
\item 全部 AI 生成的数字均由程序运行输出替换或核验：论文数值与 \texttt{results/*.json} 双向追溯；未采纳任何未经验证的 AI 断言。
\end{document}
\end{itemize}
